\chapter{पूर्ण समष्टियाँ: बेयर, आस्कोली, स्टोन--वाइरश्ट्रास}
\label{ch:b3:complete}

\omterm{def:b3:complete:complete}{पूर्णता} --- अर्थात् प्रत्येक कोशी अनुक्रम का अभिसरित होना --- वह
गुण है जो विश्लेषण को वस्तुएँ \emph{उत्पन्न} करने देता है:
संकुचनों के स्थिर बिंदु, श्रेणियों के योग, सीमा के रूप में
प्राप्त समीकरणों के हल। यह अध्याय मापीय सिद्धांत की तीन महान
अस्तित्व-मशीनें जोड़ता है। \emph{बेयर की प्रमेय} दिखाती है कि कोई
\omterm{def:b3:complete:complete}{पूर्ण} समष्टि नगण्य टुकड़ों का गणनीय सम्मिलन नहीं हो सकती, और
विशुद्ध गणनीयता जैसे तर्क से वस्तुएँ (\omterm{def:b3:topology:continuity}{संतत} पर \omterm{thm:b3:complete:nowherediff}{कहीं भी अनवकलनीय
फलन}!) प्रकट कर देती है। \emph{आर्ज़ेला--आस्कोली} $\mathcal C(K)$ के \omterm{def:b3:topology:compact}{संहत}
उपसमुच्चयों की पहचान कर देती है और विश्लेषण का संहतता-श्रमिक है
--- सप्ताहांत समस्या उससे अवकल समीकरणों के लिए पेआनो की अस्तित्व
प्रमेय सिद्ध करती है। \emph{स्टोन--वाइरश्ट्रास} दिखाती है कि
बहुपद, और उनके अतिरिक्त बहुत कुछ, $\mathcal C(K)$ में सघन हैं: आसन्नन एक
\omterm{def:b3:galois:algebraic}{बीजीय} सत्यापन बन जाता है। रास्ते में हम \omterm{thm:b3:complete:completion}{पूर्णन} बनाते हैं और
एकसमान \omterm{def:b3:topology:continuity}{संतत} प्रतिचित्रणों के लिए विस्तार प्रमेय सिद्ध करते हैं,
जो \cref{ch:b3:lp,ch:b3:hilbert,ch:b3:fouriertransform} की रोज की रोटी है।

\section{पूर्ण समष्टियाँ, पूर्णन, विस्तार}

\begin{definition}\label{def:b3:complete:complete}
कोई मापीय समष्टि \emph{पूर्ण}\index{पूर्ण मापीय समष्टि} कहलाती है
यदि उसका प्रत्येक कोशी अनुक्रम अभिसरित हो (दूसरा वर्ष: $\R^n$ पूर्ण
है; $\norm\cdot_\infty$ के साथ $\mathcal C(\intcc01)$ पूर्ण है)। किसी पूर्ण समष्टि का संवृत
उपसमुच्चय पूर्ण होता है; और किसी भी मापीय समष्टि का पूर्ण
उपसमुच्चय संवृत होता है।
\end{definition}

\begin{proof}
दोनों कथनों के लिए: संवृत $F$ का कोई कोशी अनुक्रम $X$ में
अभिसरित होता है, और उसकी सीमा, $F$ की संलग्न होने के कारण, $F$
में ही होती है; और \omterm{def:b3:complete:complete}{पूर्ण} $A$ का $X$ में अभिसरित होने वाला कोई
अनुक्रम कोशी है, अतः $A$ में अभिसरित होता है, और सीमाएँ अद्वितीय
होती हैं।
\end{proof}

\begin{theorem}[एकसमान संतत प्रतिचित्रणों का विस्तार]
\label{thm:b3:complete:extension}
मान लीजिए $D \subseteq X$ सघन है, $Y$ \omterm{def:b3:complete:complete}{पूर्ण} है, और $f \colon D \to
Y$ एकसमान \omterm{def:b3:topology:continuity}{संतत}। तब
$f$ \emph{अद्वितीय} रूप से किसी \omterm{def:b3:topology:continuity}{संतत} $\bar f \colon X \to Y$ तक विस्तारित हो जाता
है, और $\bar f$ एकसमान \omterm{def:b3:topology:continuity}{संतत} होता है।\index{विस्तार प्रमेय (एकसमान
सांतत्य)}
\end{theorem}

\begin{proof}
अद्वितीयता: दो \omterm{def:b3:topology:continuity}{संतत} विस्तार सघन $D$ पर सहमत होते हैं, अतः सर्वत्र
(सहमति समुच्चय $\{g = h\}$ संवृत है: $x\mapsto(g(x), h(x))$ के अंतर्गत संवृत विकर्ण का
पूर्वप्रतिबिंब)। अस्तित्व: $x \in X$ के लिए $d_n \to x$, $d_n \in D$ चुनिए।
अनुक्रम $(f(d_n))$ कोशी है: $\varepsilon$ दिया हो तो एकसमान \omterm{def:b3:topology:continuity}{सांतत्य} $d(u,v) < \delta \Rightarrow
d(f(u), f(v)) < \varepsilon$ वाला
$\delta$ दे देता है, और $(d_n)$ कोशी है। $\bar f(x) = \lim f(d_n)$ परिभाषित कीजिए; सीमा चुने
गए अनुक्रम पर निर्भर नहीं करती (दो अनुक्रमों को गूँथ दीजिए)।
$\bar f$ $f$ का विस्तार है (अचर अनुक्रम) और \omterm{def:b3:topology:continuity}{सांतत्य} मापांक विरासत में
पाता है: यदि $d(x, x') < \delta$ हो, तो दोनों को दूरी $< \frac{\delta - d(x,x')}2$ पर $D$ के बिंदुओं से
आसन्न करने पर सीमा में $d(\bar f(x), \bar
f(x')) \leq \varepsilon$ मिलता है --- अतः $\bar f$ एकसमान \omterm{def:b3:topology:continuity}{संतत} है।
\end{proof}

\begin{theorem}[पूर्णन]\label{thm:b3:complete:completion}
प्रत्येक मापीय समष्टि $X$ किसी \omterm{def:b3:complete:complete}{पूर्ण मापीय समष्टि} $\hat X$ के सघन
उपसमुच्चय के रूप में समदूरिक ढंग से अंतःस्थापित हो जाती है, और
$X$ को बिंदुशः स्थिर रखने वाली समदूरीयता तक वह अद्वितीय है:
यही उसका \emph{पूर्णन}\index{मापीय समष्टि का पूर्णन} है।
\end{theorem}

\begin{proof}
\emph{अस्तित्व।} मान लीजिए $\mathcal X$ $X$ के कोशी अनुक्रमों का
समुच्चय है, जिस पर छद्म-दूरी
\[
D\bigl((x_n), (y_n)\bigr) = \lim_n d(x_n, y_n),
\]
है; सीमा विद्यमान है क्योंकि $\abs{d(x_n, y_n) - d(x_m, y_m)}
\leq d(x_n, x_m) + d(y_n, y_m)$ वास्तविक अनुक्रम को कोशी बना
देता है। $\hat X = \mathcal X/{\sim}$ लीजिए, जिसमें $D$-दूरी $0$ वाले अनुक्रमों की
पहचान कर दी गई है; $D$ उतरकर एक दूरी बन जाती है। $X$ को अचर
अनुक्रमों द्वारा अंतःस्थापित कीजिए: यह समदूरीयता है और उसका
प्रतिबिंब सघन है (किसी कोशी अनुक्रम का $D$-आसन्नन उसके अपने ही
पदों पर बने अचरों से हो जाता है: कोशीपन से $k \to \infty$ पर $D\bigl((x_n), (x_k)_{\rm
const}\bigr) = \lim_n d(x_n, x_k) \to 0$)। $\hat X$
की \omterm{def:b3:complete:complete}{पूर्णता}: मान लीजिए $(\xi^k)$ $\hat X$ में कोशी है; सघनता से $D(\xi^k, x_k)
\leq 2^{-k}$ वाले
$x_k \in X$ चुनिए; तब $(x_k)$ $X$ में कोशी है ($\xi$ से होकर त्रिभुज
असमिका), जो एक बिंदु $\xi \in \hat
X$ परिभाषित करता है, और $D(\xi^k, \xi) \leq 2^{-k} + D(x_k, \xi) \to 0$ (अचर $x_k$ से
$(x_j)_j$ के वर्ग की दूरी $\lim_j d(x_k, x_j)$ है, जो बड़े $k$ के लिए छोटी है)।

\emph{अद्वितीयता}: दो \omterm{thm:b3:complete:completion}{पूर्णन} $\hat X_1, \hat X_2$ $X$ को सघन रूप से अंतर्विष्ट
करते हैं; $X$ का तत्समक, जो एक समदूरीयता है, एकसमान \omterm{def:b3:topology:continuity}{संतत} है,
अतः $\hat X_1 \to \hat X_2$ तक विस्तारित हो जाता है (\cref{thm:b3:complete:extension}), और \omterm{def:b3:topology:interior}{सघन समुच्चय} पर
समदूरीयता होने के कारण सर्वत्र समदूरीयता रहता है; सममित रूप से
दूसरी दिशा में भी, और संयोजन सघन $X$ को स्थिर रखते हैं: अतः वे
तत्समक हैं।
\end{proof}

\begin{theorem}[बानाख स्थिर बिंदु]\label{thm:b3:complete:banach}
मान लीजिए $X$ \omterm{def:b3:complete:complete}{पूर्ण} और अरिक्त है और $f \colon X \to X$ एक \emph{संकुचन}:
$k < 1$ के साथ $d(f(x), f(y)) \leq k\,d(x,y)$। तब $f$ का एक अद्वितीय स्थिर बिंदु $x^*$ होता
है, और प्रत्येक \omterm{def:b3:groups:action}{कक्षा} उस पर अभिसरित होती है, और स्पष्ट दर
$d(x_n, x^*) \leq
\frac{k^n}{1-k}\,d(x_1, x_0)$ के साथ।\index{बानाख स्थिर बिंदु प्रमेय}
\end{theorem}

\begin{proof}
(इसे दूसरे वर्ष में सिद्ध किया जा चुका है; स्वयंपूर्णता के लिए हम
दो पंक्ति का तर्क फिर दर्ज कर देते हैं।) \omterm{def:b3:groups:action}{कक्षा} $x_{n+1} = f(x_n)$ के लिए
$d(x_{n+1},
x_n) \leq k^nd(x_1, x_0)$, अतः वह कोशी है (गुणोत्तर श्रेणी); उसकी सीमा $x^*$ स्थिर है
($f$ का \omterm{def:b3:topology:continuity}{सांतत्य}), और अद्वितीय है क्योंकि दो स्थिर बिंदु $d \leq k\,d$
को संतुष्ट करते हैं। दर: गुणोत्तर पुच्छ का योग लीजिए।
\end{proof}

\begin{example}[तत्समक का विक्षोभ]
\label{ex:b3:complete:perturbident}
मान लीजिए $g \colon \R^d \to \R^d$ $k < 1$ के साथ $k$-लिप्शिट्ज़ है। तब $\varphi = \mathrm{id} + g$
\emph{$\R^d$ की $\R^d$ पर \omterm{def:b3:topology:continuity}{समस्थितिकता}} है। एकैकीपन, एक मात्रात्मक
मापांक सहित:
\[
\norm{\varphi(x) - \varphi(y)} \geq \norm{x - y} -
\norm{g(x) - g(y)} \geq (1 - k)\norm{x - y} .
\]
आच्छादकता स्थिर बिंदु प्रमेय है: $\varphi(x) = y$ हल करने का अर्थ है $x = y - g(x)$,
और $x \mapsto y -
g(x)$ \omterm{def:b3:complete:complete}{पूर्ण} $\R^d$ का $k$-संकुचन है --- अतः प्रत्येक $y$ के
लिए एक अद्वितीय हल $x = \psi(y)$ विद्यमान है। प्रदर्शित असमिका प्रतिलोम
$\psi$ को नियतांक $\frac1{1-k}$ के साथ लिप्शिट्ज़ बना देती है: अर्थात् एक
\omterm{def:b3:topology:continuity}{समस्थितिकता}, जिसके दोनों मापांकों पर स्पष्ट परिबंध हैं। यह
भोला-सा दिखने वाला कथन प्रतिलोम फलन प्रमेय (\cref{ch:b3:submanifolds}) के भीतर का
इंजन है: जिस बिंदु के निकट $Df$ व्युत्क्रमणीय है, वहाँ $f$
\emph{वस्तुतः} एक व्युत्क्रमणीय रैखिक प्रतिचित्रण और एक छोटा
लिप्शिट्ज़ विक्षोभ है, और शेष काम आज का उदाहरण कर देता है। यह
सांख्यिक दृढ़ता को भी माप देता है: प्रतिलोम के अंतराल से कम
विक्षुब्ध कोई निकाय हल्य बना रहता है, और हल विक्षोभ के अधिकतम
$\frac{1}{1-k}$ गुना ही खिसकता है।
\end{example}

\section{बेयर की प्रमेय}

\begin{theorem}[बेयर]\label{thm:b3:complete:baire}
किसी \omterm{def:b3:complete:complete}{पूर्ण मापीय समष्टि} में सघन विवृत समुच्चयों का गणनीय
प्रतिच्छेदन सघन होता है। समतुल्य रूप से: यदि $X = \bigcup_{n}F_n$ हो जहाँ
प्रत्येक $F_n$ संवृत है, तो किसी $F_n$ का \omterm{def:b3:topology:interior}{अंतःभाग} अरिक्त
होता है।\index{बेयर की प्रमेय}
\end{theorem}

\begin{proof}
मान लीजिए $(U_n)$ सघन विवृत हैं और $B_0 = B(x_0, r_0)$ कोई विवृत गोला; हम $B_0$
में $\bigcap U_n$ का कोई बिंदु ढूँढ़ते हैं। आगमन से: $U_{n}$ के सघन और
विवृत होने से वह विवृत गोले $B_{n-1}$ से किसी \omterm{def:b3:topology:topology}{विवृत समुच्चय} में
मिलता है, जो $0 < r_n \leq r_{n-1}/2$ और $\bar B(x_n, r_n) \subseteq
B_{n-1}\cap U_n$ वाले किसी संवृत गोले $\bar B(x_n, r_n)$ को
अंतर्विष्ट करता है। \omterm{ex:b3:groups:actions}{केंद्र} एक कोशी अनुक्रम बनाते हैं ($m \geq n$ के
लिए $x_m \in
B_n$, त्रिज्याएँ $\to 0$); सीमा $x$ प्रत्येक $\bar B(x_n, r_n)$ में होती
है (संवृतता), अतः प्रत्येक $U_n$ में और $B_0$ में भी। दूसरे रूप
के लिए: यदि किसी $F_n$ का \omterm{def:b3:topology:interior}{अंतःभाग} न हो, तो $U_n = X \setminus F_n$ विवृत और सघन
हैं, और $\bigcap U_n$ का कोई बिंदु $\bigcup F_n = X$ से बच निकलता है: निरर्थक।
\end{proof}

\begin{remark}\label{rem:b3:complete:meagre}
शब्दावली: कोई समुच्चय \emph{कहीं सघन नहीं} कहलाता है यदि उसकी
\omterm{def:b3:topology:interior}{संवृति} का \omterm{def:b3:topology:interior}{अंतःभाग} रिक्त हो, और \emph{क्षीण}\index{क्षीण समुच्चय}
(प्रथम श्रेणी का) यदि वह कहीं सघन नहीं समुच्चयों का गणनीय
सम्मिलन हो। बेयर: \emph{कोई \omterm{def:b3:complete:complete}{पूर्ण मापीय समष्टि} अपने आप में \omterm{rem:b3:complete:meagre}{क्षीण}
नहीं होती}, और किसी \omterm{rem:b3:complete:meagre}{क्षीण समुच्चय} का पूरक सघन होता है। ``\omterm{rem:b3:complete:meagre}{क्षीण}''
छोटेपन की ऐसी धारणा है जो माप के लंबवत है (\cref{ch:b3:measure} पूरे माप वाले
\omterm{rem:b3:complete:meagre}{क्षीण समुच्चय} बना देगा), और बेयर के तर्क \emph{प्रचुरता से
अस्तित्व} सिद्ध करते हैं: गुण P न रखने वाली एक वस्तु प्रस्तुत
करने के लिए दिखाइए कि P वाली वस्तुएँ एक \omterm{rem:b3:complete:meagre}{क्षीण समुच्चय} बनाती हैं।
\end{remark}

\begin{corollary}\label{cor:b3:complete:bairecors}
(क) $\R$ अगणनीय है। (ख) $\Q$ $\R$ के विवृत उपसमुच्चयों का
गणनीय प्रतिच्छेदन नहीं है, और विविक्त बिंदुओं के बिना कोई अरिक्त
\omterm{def:b3:complete:complete}{पूर्ण मापीय समष्टि} अगणनीय होती है।
\end{corollary}

\begin{proof}
(क) किसी गणनीय समुच्चय पर $\R = \bigcup_{x}\{x\}$ से किसी एकल-अवयवी समुच्चय का
\omterm{def:b3:topology:interior}{अंतःभाग} बन जाता। (ख) यदि विवृत $V_n$ (जो अनिवार्यतः सघन हैं,
क्योंकि $\supseteq \Q$) के साथ $\Q = \bigcap_n V_n$ हो, तो समुच्चय $V_n$ और पूरक $\R\setminus\{q\}$,
$q \in \Q$, सघन विवृत समुच्चयों का ऐसा गणनीय कुल बनाते हैं जिसका
प्रतिच्छेदन रिक्त है --- जो बेयर का खंडन है। यदि $X$ विविक्त
बिंदुओं के बिना \omterm{def:b3:complete:complete}{पूर्ण} और गणनीय हो, तो $X = \bigcup_{x \in X}\{x\}$ उसे रिक्त \omterm{def:b3:topology:interior}{अंतःभाग}
वाले संवृत समुच्चयों के गणनीय सम्मिलन के रूप में प्रस्तुत कर देता
है (कोई विविक्त बिंदु नहीं): फिर बेयर।
\end{proof}

\begin{theorem}[वाइरश्ट्रास के राक्षस विद्यमान हैं]
\label{thm:b3:complete:nowherediff}
$\intcc01$ पर ऐसे \omterm{def:b3:topology:continuity}{संतत} फलन विद्यमान हैं जो किसी भी बिंदु पर अवकलनीय
नहीं हैं। वस्तुतः, ऐसे $f \in \mathcal C(\intcc01)$ का समुच्चय जिनका एक भी बिंदु पर
(परिमित) अवकलज हो, $(\mathcal C(\intcc01), \norm\cdot_\infty)$ में \omterm{rem:b3:complete:meagre}{क्षीण} है।\index{कहीं भी अनवकलनीय
फलन}
\end{theorem}

\begin{proof}
$n \geq 1$ के लिए मान लीजिए
\[
F_n = \Bigl\{f : \exists x \in \intcc01,\
\forall h \neq 0 \text{ जिसके लिए } x + h \in \intcc01,\
\abs{f(x+h) - f(x)} \leq n\abs h \Bigr\}.
\]
यदि $f$ $x$ पर अवकलनीय हो, तो किसी $n$ के लिए $f \in F_n$:
भागफल $\abs{f(x+h)-f(x)}/\abs h$ $\abs h
\leq \delta$ के लिए परिबद्ध है (अवकलनीयता: वह $\abs{f'(x)}$ की ओर
जाता है) और $\abs h \geq \delta$ के लिए $2\norm f_\infty/\delta$ से परिबद्ध। अतः $\bigcup F_n$ समस्त
कहीं-न-कहीं अवकलनीय फलनों को अंतर्विष्ट करता है, और केवल यह
दिखाना पर्याप्त है कि प्रत्येक $F_n$ संवृत है और उसका \omterm{def:b3:topology:interior}{अंतःभाग}
रिक्त है।

\emph{संवृत}: मान लीजिए एकसमान रूप से $f_k \to f$, और साक्षियों $x_k \to x$
के साथ $f_k \in F_n$ (संहतता, निष्कर्षण के बाद)। $x + h \in \intcc01$ वाले $h$ के लिए:
$x_k + h_k
\in \intcc01$ वाला $h_k \to h$ चुनिए (उदाहरणार्थ काटा हुआ $h_k = h + x - x_k$); तब एकसमान
अभिसरण और संगत बिंदुओं पर $f$ के \omterm{def:b3:topology:continuity}{सांतत्य} का उपयोग करते हुए
$\abs{f(x + h) - f(x)} = \lim \abs{f_k(x_k + h_k) - f_k(x_k)}
\leq \lim n\abs{h_k} = n\abs h$: अतः $f \in F_n$।

\emph{रिक्त \omterm{def:b3:topology:interior}{अंतःभाग}}: $f \in F_n$ और $\varepsilon > 0$ दिए हों, तो हम $\norm{g - f}_\infty \leq \varepsilon$ और
$g
\notin F_n$ वाला $g$ ढूँढ़ते हैं। पहले $f$ का $\varepsilon/2$ के भीतर किसी
खंडशः आफ़ीन $\varphi$ से आसन्नन कीजिए (एकसमान \omterm{def:b3:topology:continuity}{सांतत्य}: किसी सूक्ष्म
जालिका पर अंतर्वेशन कीजिए), जिसकी ढालें किसी $M$ से परिबद्ध
हों। अब एक छोटा आरादंत जोड़िए: $g = \varphi + \frac\varepsilon2\,s_N$, जहाँ $s_N(x)$ आयाम $1$ और
ढाल $\pm 2N$ वाला $\frac1N$-आवर्ती टेढ़ा-मेढ़ा फलन है। प्रत्येक $x$
पर, किसी एक ओर मनमाना छोटा $h$ ऐसा है कि आरादंत $[x, x+h]$ पर ढाल
$\pm 2N\cdot\frac\varepsilon2 =
\pm\varepsilon N$ का योगदान देता है: अतः बड़े $N$ के लिए $g$ का अंतर
भागफल $\varepsilon N - M > n$ से अधिक हो जाता है। इसलिए $g \notin
F_n$, और यह $f$ से किसी
भी एकसमान दूरी $\leq \varepsilon$ पर। निष्कर्ष: $\bigcup F_n$ \omterm{rem:b3:complete:meagre}{क्षीण} है; और बेयर से
उसका पूरक --- जो कहीं भी अनवकलनीय फलनों से बना है --- $\mathcal C(\intcc01)$ में
सघन है: ऐसे फलन \emph{प्रचुरता में} विद्यमान हैं।
\end{proof}

\section{आर्ज़ेला--आस्कोली}

आगे सर्वत्र $K$ एक \emph{\omterm{def:b3:topology:compact}{संहत}} मापीय समष्टि है और $\norm f_\infty = \sup_K
\norm{f(x)}$ के साथ
$\mathcal
C(K) = \mathcal C(K, \R^d)$: एक \omterm{def:b3:complete:complete}{पूर्ण} समष्टि (\omterm{def:b3:topology:continuity}{संतत} फलनों की एकसमान सीमाएँ \omterm{def:b3:topology:continuity}{संतत} होती
हैं --- दूसरा वर्ष)।

\begin{definition}\label{def:b3:complete:equicontinuous}
कोई कुल $\mathcal F \subseteq \mathcal C(K)$ \emph{समसंतत}\index{समसांतत्य} कहलाता है यदि
प्रत्येक $\varepsilon > 0$ के लिए ऐसा $\delta > 0$ हो कि
\[
d(x, y) < \delta \implies \norm{f(x) - f(y)} < \varepsilon
\quad\text{\emph{सभी} } f \in \mathcal F \text{ के लिए}
\]
(पूरे कुल के लिए एक ही $\delta$ --- उदाहरणार्थ उभयनिष्ठ लिप्शिट्ज़
नियतांक वाला कोई कुल, अथवा उभयनिष्ठ होल्डर मापांक वाला), और
\emph{बिंदुशः परिबद्ध} यदि प्रत्येक $x$ के लिए $\sup_{f}\norm{f(x)} < \infty$ हो।
\end{definition}

\begin{theorem}[आर्ज़ेला--आस्कोली]\label{thm:b3:complete:ascoli}
कोई उपसमुच्चय $\mathcal F \subseteq \mathcal C(K)$ सापेक्षतः \omterm{def:b3:topology:compact}{संहत} है (उसकी \omterm{def:b3:topology:interior}{संवृति} \omterm{def:b3:topology:compact}{संहत} है) तभी
और केवल तभी जब वह \omterm{def:b3:complete:equicontinuous}{समसंतत} और बिंदुशः परिबद्ध हो। विशेष रूप से,
प्रत्येक \omterm{def:b3:complete:equicontinuous}{समसंतत}, बिंदुशः परिबद्ध \emph{अनुक्रम} का कोई एकसमान
अभिसारी उपानुक्रम होता है।\index{आर्ज़ेला--आस्कोली प्रमेय}
\end{theorem}

\begin{proof}
($\Leftarrow$) मान लीजिए $(f_k)$ $\mathcal F$ का कोई अनुक्रम है। \omterm{def:b3:topology:compact}{संहत} मापीय
$K$ \emph{पृथक्करणीय} है: प्रत्येक $n$ के लिए त्रिज्या $\frac1n$
के परिमित संख्या के गोले $K$ को आच्छादित कर देते हैं (\cref{thm:b3:topology:metriccompact});
उनके \omterm{ex:b3:groups:actions}{केंद्र} एक गणनीय \omterm{def:b3:topology:interior}{सघन समुच्चय} $D = \{x_1, x_2, \dots\}$ बनाते हैं। बिंदुशः
परिबद्धता और बोलत्सानो--वाइरश्ट्रास से क्रमशः $x_1$ पर, फिर
$x_2$ पर भी, अभिसरित होने वाले उपानुक्रम निकालिए, इत्यादि, और
\emph{विकर्ण} उपानुक्रम $(g_j)$ लीजिए: वह $D$ के प्रत्येक बिंदु
पर अभिसरित होता है। \omterm{def:b3:complete:equicontinuous}{समसांतत्य} इसे एकसमान कोशीपन तक उठा देता है:
$\varepsilon$ दिया हो तो परिभाषा के अनुसार $\delta$ लीजिए, $K$ को $x_i \in D$
($i \leq m$) वाले परिमित संख्या के गोलों $B(x_i, \delta)$ से आच्छादित कीजिए, और
इतना बड़ा $J$ चुनिए कि $j, j' \geq
J$, $i \leq m$ के लिए $\norm{g_j(x_i) - g_{j'}(x_i)} < \varepsilon$ हो। मनमाने
$x \in B(x_i, \delta)$ के लिए:
\[
\norm{g_j(x) - g_{j'}(x)}
\leq \norm{g_j(x) - g_j(x_i)} + \norm{g_j(x_i) - g_{j'}(x_i)}
+ \norm{g_{j'}(x_i) - g_{j'}(x)} < 3\varepsilon .
\]
अतः $(g_j)$ एकसमान कोशी है और \omterm{def:b3:complete:complete}{पूर्ण} $\mathcal C(K)$ में अभिसरित होता है। इस
प्रकार $\mathcal F$ के प्रत्येक अनुक्रम का कोई अभिसारी उपानुक्रम है:
अर्थात् $\bar{\mathcal F}$ (अनुक्रमीय रूप से, और इसलिए \cref{thm:b3:topology:metriccompact} से) \omterm{def:b3:topology:compact}{संहत} है।

($\Rightarrow$) यदि $\bar{\mathcal F}$ \omterm{def:b3:topology:compact}{संहत} हो: बिंदुशः परिबद्धता स्पष्ट है
(मान-निर्धारण \omterm{def:b3:topology:continuity}{संतत} है)। \omterm{def:b3:complete:equicontinuous}{समसांतत्य} के लिए $\mathcal F$ को $\mathcal C(K)$ के परिमित
संख्या के गोलों $B(f_i, \varepsilon)$ से आच्छादित कीजिए; प्रत्येक $f_i$ एकसमान
\omterm{def:b3:topology:continuity}{संतत} है (हाइने, \cref{cor:b3:topology:heineborel}), जिससे $i \leq m$ के लिए एक उभयनिष्ठ $\delta$
मिलता है; तब $f \in B(f_i, \varepsilon)$ और $d(x,y)
< \delta$ के लिए $\norm{f(x)-f(y)} \leq 2\varepsilon +
\norm{f_i(x)-f_i(y)} < 3\varepsilon$।
\end{proof}

\begin{example}\label{ex:b3:complete:ascoliexamples}
$\mathcal C(\intcc01)$ का संवृत इकाई गोला \omterm{def:b3:topology:compact}{संहत} \emph{नहीं} है ($f_n(x) = x^n$ का कोई
एकसमान अभिसारी उपानुक्रम नहीं है: बिंदुशः सीमा असंतत है), और
वस्तुतः $(x^n)$ $1$ पर \omterm{def:b3:complete:equicontinuous}{समसंतत} नहीं है। इसके विपरीत $\{f : \norm
f_\infty \leq 1,\ \operatorname{Lip}(f) \leq 1\}$ \omterm{def:b3:topology:compact}{संहत}
है: परिबद्ध, $1$-लिप्शिट्ज़-समसंतत, और संवृत। आस्कोली समझा देती
है कि अनंत विमा में संहतता \emph{क्यों} विफल हो जाती है (रीस,
दूसरा वर्ष) और उसे बहाल करने के लिए क्या जोड़ना पड़ता है: एक
एकसमान \omterm{def:b3:topology:continuity}{सांतत्य} मापांक।
\end{example}

\section{स्टोन--वाइरश्ट्रास}

\begin{lemma}[दीनी]\label{lem:b3:complete:dini}
मान लीजिए $K$ \omterm{def:b3:topology:compact}{संहत} है और $(f_n)$ \omterm{def:b3:topology:continuity}{संतत} वास्तविक फलनों का ऐसा
\emph{एकदिष्ट} अनुक्रम है जो किसी \omterm{def:b3:topology:continuity}{संतत} $f$ पर \emph{बिंदुशः}
अभिसरित होता है। तब अभिसरण एकसमान है।\index{दीनी की प्रमेय}
\end{lemma}

\begin{proof}
मान लीजिए $f_n \uparrow f$; तब $g_n = f - f_n \downarrow 0$ \omterm{def:b3:topology:continuity}{संतत} है। $\varepsilon$ दिया हो, तो \omterm{def:b3:topology:topology}{विवृत
समुच्चय} $U_n = \{g_n <
\varepsilon\}$ बढ़ते जाते हैं और $K$ को आच्छादित करते हैं
(बिंदुशः अभिसरण); कोई परिमित उपआवरण निकालिए: किसी $n_0$ के लिए
$K = U_{n_0}$ (बढ़ता हुआ कुल), अर्थात् $n \geq n_0$ के लिए सर्वत्र $0 \leq g_n < \varepsilon$।
\end{proof}

\begin{lemma}\label{lem:b3:complete:sqrt}
ऐसे \emph{बहुपदों} $u_n$ का एक अनुक्रम है जिसके लिए $\intcc01$ पर
एकसमान रूप से $u_n(t) \to
\sqrt t$ हो।
\end{lemma}

\begin{proof}
$u_0 = 0$, $u_{n+1}(t) = u_n(t) + \frac12\bigl(t -
u_n(t)^2\bigr)$ परिभाषित कीजिए: ये बहुपद हैं। $\intcc01$ पर आगमन से
$0 \leq u_n(t) \leq
\sqrt t$: $n$ के लिए इसे मानकर,
\[
\sqrt t - u_{n+1}(t) = \bigl(\sqrt t - u_n(t)\bigr)
\Bigl(1 - \tfrac12\bigl(\sqrt t + u_n(t)\bigr)\Bigr) \geq 0,
\]
क्योंकि $\sqrt t + u_n \leq 2$; और $u_{n+1} \geq u_n \geq 0$। अतः $(u_n(t))$ अनवरोही है और $\sqrt t$ से
परिबद्ध: वह बिंदुशः अभिसरित होता है, और सीमा $\ell(t)$ $\ell =
\ell + \frac12(t - \ell^2)$ को
संतुष्ट करती है: अर्थात् $\ell(t) = \sqrt t$, जो \omterm{def:b3:topology:continuity}{संतत} है। दीनी (\cref{lem:b3:complete:dini}) इसे
एकसमान तक उठा देती है।
\end{proof}

\begin{theorem}[स्टोन--वाइरश्ट्रास, वास्तविक रूप]
\label{thm:b3:complete:stoneweierstrass}
मान लीजिए $K$ एक \omterm{def:b3:topology:compact}{संहत} (\omterm{def:b3:topology:hausdorff}{हाउसडॉर्फ}) समष्टि है और $\mathcal A \subseteq
\mathcal C(K, \R)$ ऐसा
\emph{उपबीजगणित} (योग, गुणनफल और अदिश गुणजों के अंतर्गत स्थिर) है
जो \emph{अचरों को अंतर्विष्ट करता है} और \emph{बिंदुओं को पृथक
करता है} ($x \neq y$ के लिए किसी $f \in
\mathcal A$ के लिए $f(x) \neq f(y)$)। तब $\mathcal A$ $(\mathcal C(K,\R), \norm\cdot_\infty)$
में सघन है।\index{स्टोन--वाइरश्ट्रास प्रमेय}
\end{theorem}

\begin{proof}
मान लीजिए $\bar{\mathcal A}$ उसकी \omterm{def:b3:topology:interior}{संवृति} है, जो पुनः एक बीजगणित है (परिबद्ध
समुच्चयों पर एकसमान सीमाओं के गुणनफल अभिसरित होते हैं)।

\emph{चरण 1: $\bar{\mathcal A}$ एक जालक है}, अर्थात् $\max$ और $\min$ के अंतर्गत
स्थिर। चूँकि $\max(f,g) = \frac{f + g +
\abs{f-g}}2$ और वैसे ही $\min$, इतना ही पर्याप्त है कि
$f \in
\bar{\mathcal A} \Rightarrow \abs f \in \bar{\mathcal A}$: $M = \norm f_\infty > 0$ के साथ $\abs f = M\sqrt{(f/M)^2}$, और \cref{lem:b3:complete:sqrt} से ऐसे बहुपद $u_n$ मिलते हैं
जिनके लिए एकसमान रूप से $u_n\bigl((f/M)^2\bigr) \to \abs f/M$; और बीजगणित के सदस्यों के बहुपद
(अचर पद सहित: अचर वहाँ हैं ही) $\bar{\mathcal A}$ में ही रहते हैं।

\emph{चरण 2: दो-बिंदु अंतर्वेशन।} $x \neq y$ और $a,
b \in \R$ के लिए किसी
$g \in \mathcal A$ के लिए $g(x) = a$, $g(y) = b$: $x, y$ को पृथक करने वाला $h$ लीजिए
और $g = a + (b -
a)\frac{h - h(x)}{h(y) - h(x)}$ रखिए।

\emph{चरण 3.} मान लीजिए $f \in \mathcal C(K)$, $\varepsilon > 0$। प्रत्येक युग्म $x, y$ के
लिए $g_{x,y}(x) = f(x)$, $g_{x,y}(y) = f(y)$ वाला $g_{x,y} \in \mathcal A$ चुनिए (चरण 2; $x = y$ के लिए अचर फलन
$g_{x,x} = f(x)$ लीजिए)। $x$ स्थिर कीजिए: प्रत्येक $y$ के लिए \omterm{def:b3:topology:topology}{विवृत
समुच्चय} $V_y = \{g_{x,y} < f +
\varepsilon\}$ $y$ को अंतर्विष्ट करता है; संहतता से $K = \bigcup V_{y_j}$ वाले
$y_1, \dots,
y_m$ निकल आते हैं, और $h_x = \min_j g_{x, y_j}
\in \bar{\mathcal A}$ (चरण 1) सर्वत्र $h_x < f +
\varepsilon$ को संतुष्ट
करता है, तथा $h_x(x) = f(x)$। अब $x$ को बदलिए: $W_x =
\{h_x > f - \varepsilon\}$ विवृत है और $x$
को अंतर्विष्ट करता है; $K$ को आच्छादित करने वाले $x_1,
\dots, x_l$
निकालिए, और $h = \max_i h_{x_i} \in
\bar{\mathcal A}$ $f - \varepsilon < h < f +
\varepsilon$ को संतुष्ट करता है: अतः $\norm{f - h}_\infty \leq \varepsilon$। इसलिए
$f \in \bar{\mathcal A}$।
\end{proof}

\begin{corollary}\label{cor:b3:complete:weierstrass}
(क) (\emph{वाइरश्ट्रास}) बहुपद $\mathcal
C(\intcc ab, \R)$ में सघन हैं; और \omterm{def:b3:topology:compact}{संहत} $K \subseteq \R^n$
के लिए $n$ चरों के बहुपद $\mathcal C(K, \R)$ में सघन हैं।
(ख) (\emph{सम्मिश्र रूप}) यदि $\mathcal A \subseteq \mathcal
C(K, \C)$ ऐसा उपबीजगणित हो जो अचरों को
अंतर्विष्ट करता हो, बिंदुओं को पृथक करता हो, और \emph{संयुग्मन के
अंतर्गत स्थिर} हो, तो वह सघन है।
(ग) (\emph{त्रिकोणमितीय रूप}) त्रिकोणमितीय बहुपद $\sum_{\abs n \leq N}c_n\eu^{\iu n t}$ $\norm\cdot_\infty$ के
साथ \omterm{def:b3:topology:continuity}{संतत} $2\pi$-आवर्ती फलनों की समष्टि में सघन
हैं।\index{वाइरश्ट्रास आसन्नन प्रमेय}
\end{corollary}

\begin{proof}
(क) बहुपद अचरों सहित एक बीजगणित बनाते हैं; और निर्देशांक फलन
$\R^n$ के बिंदुओं को पृथक करते हैं।
(ख) $\mathcal A$ के सदस्यों के वास्तविक और काल्पनिक भाग $\frac{f + \bar f}2$, $\frac{f - \bar f}{2\iu}$
अचरों सहित एक वास्तविक बीजगणित $\mathcal A_\R \subseteq \mathcal C(K, \R)$ बनाते हैं; वह बिंदुओं को
पृथक करता है ($f(x) \ne f(y)$ से $\operatorname{Re}f$ या $\operatorname{Im}f$ का पृथक करना अनिवार्य हो
जाता है)। वास्तविक प्रमेय लगाइए और पुनः जोड़ लीजिए।
(ग) $2\pi$-आवर्ती फलनों को $\mathcal C(S^1, \C)$ के रूप में देखिए ($S^1 = \R/2\pi\Z$, जो \omterm{def:b3:topology:compact}{संहत}
है: \cref{exo:b3:topology:5}); $\eu^{\iu t}, \eu^{-\iu t}$ और अचरों से जनित बीजगणित संयुग्मन के अंतर्गत
स्थिर है और वृत्त के बिंदुओं को पृथक करता है ($\eu^{\iu t}$ उस पर एकैकी
है)। अब (ख) लगाइए।
\end{proof}

\begin{remark}\label{rem:b3:complete:swremark}
त्रिकोणमितीय रूप दूसरे वर्ष के फूरिये अध्याय में छूटे अंतर की
मरम्मत कर देता है, और उसका विशाल सामान्यीकरण भी: $(\mathcal C(S^1), \norm\cdot_2)$ में
त्रिकोणमितीय बहुपदों की सघनता उससे और भी सहजता से निकल आती है
(मानकीकरण नियतांक तक $\norm\cdot_2 \leq \norm\cdot_\infty$), जिससे फूरिये निकाय \cref{ch:b3:hilbert} में एक
प्रसामान्य लांबिक \emph{आधार} बन जाएगा और अंततः पार्सेवाल पूरी
व्यापकता में सिद्ध हो जाएगा।
\end{remark}

\section{अभ्यास}

\begin{exercise}[$\star$]\label{exo:b3:complete:1}
(क) दर्शाइए कि $\norm f_1 = \int_0^1\abs f$ के साथ $\mathcal C(\intcc01, \R)$ \omterm{def:b3:complete:complete}{पूर्ण} \emph{नहीं} है: $\intcc0{\frac12}$ पर
$0$ के बराबर, $\intcc{\frac12 + \frac1n}1$ पर $1$ के बराबर और बीच में आफ़ीन फलन
$f_n$ $\norm\cdot_1$-कोशी हैं पर उनकी कोई \omterm{def:b3:topology:continuity}{संतत} सीमा नहीं है।
(ख) दर्शाइए कि जिस मानकित समष्टि में प्रत्येक निरपेक्षतः अभिसारी
श्रेणी अभिसरित होती है वह \omterm{def:b3:complete:complete}{पूर्ण} होती है। \emph{(किसी कोशी अनुक्रम
से $\norm{x_{n_{k+1}} -
x_{n_k}} \leq 2^{-k}$ वाला उपानुक्रम निकालिए।)}
\end{exercise}

\begin{exercise}[$\star$]\label{exo:b3:complete:2}
बेयर का उपयोग करते हुए: (क) दर्शाइए कि किसी \omterm{def:b3:complete:complete}{पूर्ण} मानकित समष्टि
का कोई गणनीय (\omterm{def:b3:galois:algebraic}{बीजीय}) आधार नहीं होता --- इससे निष्कर्ष निकालिए कि
बहुपदों की समष्टि \emph{किसी भी} मानक के लिए \omterm{def:b3:complete:complete}{पूर्ण} नहीं है;
(ख) दर्शाइए कि यदि \omterm{def:b3:topology:continuity}{संतत} $f_n \colon \R \to \R$ का कोई अनुक्रम बिंदुशः $f$ पर
अभिसरित हो, तो $f$ के \omterm{def:b3:topology:continuity}{सांतत्य} बिंदुओं का समुच्चय सघन है।
\emph{((ख) के लिए: यह स्वीकार कीजिए या सिद्ध कीजिए कि $\Omega_\delta = \{x :
\operatorname{osc}_x f < \delta\}$
विवृत है, और $\R$ को आच्छादित करने वाले संवृत समुच्चयों $F_{N} = \{x: \abs{f_n(x) - f_m(x)} \leq \delta/3\
\forall n,m \geq N\}$ का
उपयोग करके उसे सघन दिखाइए; बेयर लगाने के लिए किसी मनमाने संवृत
गोले में कार्य कीजिए।)}
\end{exercise}

\begin{exercise}[$\star\star$]\label{exo:b3:complete:3}
(क) दर्शाइए कि $f(x) = \frac12\bigl(x + \frac ax\bigr)$ ($a >
1$) $[\sqrt a, +\infty)$ का संकुचन है और उसका स्थिर
बिंदु पहचानिए --- हेरोन की विधि। $a = 2$ के लिए $x_0 = a$ से आरंभ करके
$10^{-12}$-यथार्थता के लिए पुनरावृत्तियों की संख्या आँकिए।
(ख) (केपलर समीकरण) $0 \leq e < 1$ और $m \in \R$ के लिए दर्शाइए कि $x = m + e\sin x$ का एक
अद्वितीय हल है, जो $m$ पर संततता से निर्भर करता है।
\end{exercise}

\begin{exercise}[$\star\star$]\label{exo:b3:complete:4}
(क) मान लीजिए $X$ \omterm{def:b3:complete:complete}{पूर्ण} है और $f \colon X \to X$ ऐसा है कि उसकी कोई
पुनरावृत्ति $f^p$ संकुचन है। दर्शाइए कि $f$ का एक अद्वितीय
स्थिर बिंदु है। अनुप्रयोग: $\mathcal C(\intcc0a)$ पर समाकल संकारक $T$, $Tf(x) = \int_0^x f$,
$\norm{T^p} \leq a^p/p!$ को संतुष्ट करता है --- इससे निष्कर्ष निकालिए कि प्रत्येक
$\lambda$ के लिए $u = g + \lambda Tu$ हल्य है।
(ख) (एडेलश्टाइन) मान लीजिए $K$ \emph{\omterm{def:b3:topology:compact}{संहत}} है और $f\colon K \to K$ ऐसा कि
$x \neq y$ के लिए $d(f(x), f(y)) < d(x, y)$। दर्शाइए कि $f$ का एक अद्वितीय स्थिर बिंदु
है, पर संकुचन दर खो सकती है: $X = [1, +\infty)$ (\omterm{def:b3:complete:complete}{पूर्ण}, पर \omterm{def:b3:topology:compact}{संहत} नहीं) पर
$f(x) = x
+ \frac1x$ का कोई स्थिर बिंदु नहीं है, यद्यपि वह दूरियाँ कड़ाई से घटाता
है।
\end{exercise}

\begin{exercise}[$\star\star$]\label{exo:b3:complete:5}
(क) किसी \omterm{def:b3:topology:hausdorff}{हाउसडॉर्फ समष्टि} में जाने वाले दो \omterm{def:b3:topology:continuity}{संतत प्रतिचित्रण} जो
किसी सघन उपसमुच्चय पर सहमत हों सर्वत्र सहमत होते हैं; इस अध्याय
में इसका उपयोग कहाँ हुआ?
(ख) मान लीजिए $D \subseteq X$ सघन है और $f \colon D \to Y$ \omterm{def:b3:complete:complete}{पूर्ण} $Y$ के किसी सघन
उपसमुच्चय पर एक समदूरिक एकैकी आच्छादक प्रतिचित्रण, जहाँ $X$
\omterm{def:b3:complete:complete}{पूर्ण} है। दर्शाइए कि $f$ किसी समदूरिक एकैकी आच्छादक प्रतिचित्रण
$X \to Y$ तक विस्तारित हो जाता है। इससे पूर्णनों की अद्वितीयता पुनः
निकालिए।
\end{exercise}

\begin{exercise}[$\star\star$]\label{exo:b3:complete:6}
निम्नलिखित में से कौन-से कुल $\mathcal C(\intcc01)$ में \omterm{def:b3:complete:equicontinuous}{समसंतत}, बिंदुशः परिबद्ध
और सापेक्षतः \omterm{def:b3:topology:compact}{संहत} हैं?
\[
\{x \mapsto \sin(nx)\}_{n},\qquad
\{x \mapsto x^n\}_n,
\]
\[
\Bigl\{f \in \mathcal C^1 : \norm f_\infty \leq 1,\
\norm{f'}_\infty \leq 1\Bigr\},\qquad
\{f : \operatorname{Lip}(f)\leq 1,\ f(0) = 0\}.
\]
प्रत्येक उत्तर को आस्कोली से अथवा किसी प्रतिउदाहरण अनुक्रम से
उचित ठहराइए।
\end{exercise}

\begin{exercise}[$\star\star\star$]\label{exo:b3:complete:7}
(\omterm{def:b3:topology:compact}{संहत} समाकल संकारक) मान लीजिए $k \in \mathcal C(\intcc01^2)$ और $f \in \mathcal C(\intcc01)$ के लिए $Tf(x) =
\int_0^1k(x,y)f(y)\,\dd y$।
(क) दर्शाइए कि $T$ $\mathcal C(\intcc01)$ के इकाई गोले को किसी \omterm{def:b3:complete:equicontinuous}{समसंतत}, एकसमान
परिबद्ध समुच्चय पर भेजता है; निष्कर्ष निकालिए कि $T$ एक
\emph{\omterm{def:b3:topology:compact}{संहत} संकारक} है: परिबद्ध समुच्चयों के प्रतिबिंब सापेक्षतः
\omterm{def:b3:topology:compact}{संहत} होते हैं।
(ख) इससे निष्कर्ष निकालिए कि यदि $(f_n)$ परिबद्ध हो, तो $(Tf_n)$ का
कोई एकसमान अभिसारी उपानुक्रम है, और यह कि $T$ \omterm{def:b3:topology:continuity}{संतत} प्रतिलोम
वाला एकैकी आच्छादक प्रतिचित्रण नहीं हो सकता। \emph{(इकाई गोले का
प्रतिबिंब $\mathcal C(\intcc01)$ में $0$ का \omterm{def:b3:topology:compact}{संहत} \omterm{def:b3:topology:topology}{प्रतिवेश} होता: जिसे दूसरे वर्ष
की रीस प्रमेय वर्जित करती है।)}
\end{exercise}

\begin{exercise}[$\star\star$]\label{exo:b3:complete:8}
\omterm{def:b3:topology:continuity}{संतत} $f_n \colon \intcc01 \to \R$ के लिए सिद्ध या खंडित कीजिए:
(क) बिंदुशः $f_n \downarrow 0$ $\Rightarrow$ एकसमान (दीनी --- इसे फिर से सिद्ध
कीजिए); (ख) वही, एकदिष्टता के बिना; (ग) वही, एकदिष्टता के साथ पर
$f$ असंतत होते हुए; (घ) वही, एकदिष्टता और \omterm{def:b3:topology:continuity}{संतत} सीमा के साथ,
परंतु $\intoo01$ पर।
\end{exercise}

\begin{exercise}[$\star\star$]\label{exo:b3:complete:9}
(क) (आघूर्ण निर्धारण कर देते हैं) मान लीजिए $f \in \mathcal C(\intcc01)$ ऐसा है कि सभी
$n \in \N$ के लिए $\int_0^1 f(x)\,x^n\,\dd x = 0$। दर्शाइए कि $f =
0$। \emph{($f$ का बहुपदों से
एकसमान आसन्नन कीजिए और $\int f^2$ परिकलित कीजिए।)}
(ख) दर्शाइए कि सम बहुपद $\mathcal
C(\intcc01)$ में सघन हैं पर $\mathcal C(\intcc{-1}1)$ में
\emph{नहीं}; स्टोन--वाइरश्ट्रास की कौन-सी परिकल्पना विफल होती है?
(ग) क्या केवल $x \mapsto \eu^{\iu x}$ से ($\eu^{-\iu x}$ के बिना) जनित बीजगणित $\mathcal C(S^1, \C)$ में
सघन है? \emph{($\int_0^{2\pi}f(t)\,\eu^{\iu t}\dd t$ पर विचार कीजिए।)}
\end{exercise}

\begin{exercise}[$\star\star\star$]\label{exo:b3:complete:10}
(एकसमान परिबद्धता, मापीय रूप) मान लीजिए $X$ एक \omterm{def:b3:complete:complete}{पूर्ण मापीय
समष्टि} है और $\mathcal F \subseteq \mathcal C(X, \R)$ ऐसा कुल जो \emph{बिंदुशः} परिबद्ध है: प्रत्येक
$x$ के लिए $\sup_{f \in \mathcal
F}\abs{f(x)} < \infty$। दर्शाइए कि कोई अरिक्त विवृत $U \subseteq X$ है जिस पर
$\mathcal F$ \emph{एकसमान} परिबद्ध है: $\sup_{f}\sup_U \abs f < \infty$। \emph{($F_n = \{x : \abs{f(x)} \leq n\ \forall f\}$ पर विचार
कीजिए।)} यही \cref{ch:b3:banach} में बानाख--श्टाइनहाउस के पीछे का इंजन है।
\end{exercise}

\begin{exercise}[$\star\star$]\label{exo:b3:complete:11}
($\mathcal C^1$ को अपने ही मानक की आवश्यकता है) $E = \mathcal
C^1(\intcc01, \R)$ पर $\norm f_{\mathcal C^1} = \norm
f_\infty + \norm{f'}_\infty$ पर विचार
कीजिए।
(क) दर्शाइए कि $(E, \norm\cdot_{\mathcal C^1})$ \omterm{def:b3:complete:complete}{पूर्ण} है \emph{(किसी $\mathcal C^1$-कोशी अनुक्रम के
लिए $f_n \to f$ और $f_n' \to g$ एकसमान रूप से होते हैं; $f_n(x) = f_n(0) + \int_0^xf_n'$ में सीमा लेकर
$g = f'$ की पहचान कीजिए)}।
(ख) दर्शाइए कि $(E, \norm\cdot_\infty)$ \omterm{def:b3:complete:complete}{पूर्ण} \emph{नहीं} है: $\mathcal C^1$ फलनों की ऐसी
एकसमान सीमा प्रस्तुत कीजिए जो अवकलनीय न हो (उदाहरणार्थ $\abs{x - \tfrac12}$ के
चिकने आसन्नन)।
(ग) (क), (ख) और विवृत प्रतिचित्रण संबंधी विचारों से --- अथवा
सीधे --- निष्कर्ष निकालिए कि $E$ पर कोई नियतांक $C$ $\norm{f'}_\infty \leq C\,\norm f_\infty$ को
संतुष्ट नहीं करता: इसका साक्षी एक अनुक्रम प्रस्तुत कीजिए। अवकलन
अपरिबद्ध है; यही \cref{thm:b3:complete:nowherediff} के पीछे का ढलान है।
\end{exercise}

\begin{exercise}[$\star\star\star$]\label{exo:b3:complete:12}
(क्रॉफ्ट की प्रमेयिका) मान लीजिए $f\colon\intoo0\infty\to\R$ \omterm{def:b3:topology:continuity}{संतत} है और मान लीजिए कि
प्रत्येक $x > 0$ के लिए पूर्णांक $n \to \infty$ पर $f(nx) \to 0$। दर्शाइए कि $t \to
+\infty$
पर $f(t) \to 0$। \emph{($\varepsilon > 0$ स्थिर कीजिए; समुच्चय
\[
F_N = \{x > 0 : \abs{f(nx)} \leq \varepsilon\ \ \forall n
\geq N\}
\]
संवृत हैं और $\intoo0\infty$ को आच्छादित करते हैं; किसी अंतराल $\intcc ab$ में
बेयर $N$ और कोई उपअंतराल $\intcc{a'}{b'}
\subseteq F_N$ दे देता है; तब $n(b' - a') \geq a'$ हो जाने पर
प्रसार $\bigl[na', nb'\bigr]$, $n
\geq N$, $+\infty$ के एक पूरे \omterm{def:b3:topology:topology}{प्रतिवेश} को आच्छादित कर
देते हैं।)} परिकल्पना ``प्रत्येक $x$ के लिए'' (केवल परिमेय
$x$ के लिए नहीं) का उपयोग कहाँ होता है?
\end{exercise}

\section{समस्या: पेआनो की अस्तित्व प्रमेय}

\begin{problem}[{सप्ताहांत समस्या --- लिप्शिट्ज़ के बिना $x' = f(t, x)$ के
हलों का अस्तित्व}]\label{pb:b3:complete:1}
कोशी--लिप्शिट्ज़ (दूसरा वर्ष; \cref{ch:b3:ode} में पुनः सिद्ध) $f$ से
$x$ में लिप्शिट्ज़ होने की माँग करती है। पेआनो (1890):
\emph{$f$ का \omterm{def:b3:topology:continuity}{संतत} होना ही अस्तित्व दे देता है} --- यद्यपि
अद्वितीयता नहीं। हम इसे ऑयलर बहुभुजों और आस्कोली से सिद्ध करते
हैं। व्यवस्था: $f \colon R
\to \R^d$ आयत $R = \intcc{t_0 - a}{t_0
+ a}\times \bar B(x_0, b) \subseteq \R\times\R^d$ पर \omterm{def:b3:topology:continuity}{संतत} है, $M =
\sup_R\norm f$, और
\[
T = \min\bigl(a,\ b/M\bigr)
\qquad (\text{यदि } M = 0 \text{ हो तो समस्या तुच्छ है}).
\]

\medskip
\noindent\textbf{भाग I --- ऑयलर बहुभुज।} $n \geq 1$ के लिए $[t_0, t_0 + T]$ का
$t_k = t_0 + kT/n$ ($0 \leq k
\leq n$) द्वारा उपविभाजन कीजिए और $\varphi_n$ को खंडशः आफ़ीन रूप
से परिभाषित कीजिए: $\varphi_n(t_0) = x_0$ और $[t_k, t_{k+1}]$ पर
\[
\varphi_n(t) = \varphi_n(t_k) + (t -
t_k)\,f\bigl(t_k, \varphi_n(t_k)\bigr).
\]
\begin{enumerate}
  \item आगमन से दर्शाइए कि $\varphi_n$ सुपरिभाषित है, और $[t_0, t_0 + T]$ पर
        $\norm{\varphi_n(t) - x_0} \leq M(t - t_0) \leq b$ --- अतः मान-निर्धारण बिंदु $R$ में ही रहते हैं।
        \emph{(यहीं $T \leq b/M$ प्रवेश करता है।)}
  \item दर्शाइए कि प्रत्येक $\varphi_n$ $M$-लिप्शिट्ज़ है।
  \item आर्ज़ेला--आस्कोली (\cref{thm:b3:complete:ascoli}) से निष्कर्ष निकालिए कि कोई
        उपानुक्रम $(\varphi_{n_j})$ $[t_0, t_0 +
        T]$ पर एकसमान रूप से किसी $\varphi$ पर
        अभिसरित होता है, जो स्वयं $\varphi(t_0) = x_0$ के साथ $M$-लिप्शिट्ज़
        है।
\end{enumerate}

\medskip
\noindent\textbf{भाग II --- सीमा समीकरण को हल कर देती है।}
अजालिका बिंदुओं पर \emph{दोष} $\Delta_n(t) = \varphi_n'(t) -
f\bigl(t, \varphi_n(t)\bigr)$ परिभाषित कीजिए।
\begin{enumerate}[resume]
  \item दर्शाइए कि $f$ $R$ पर एकसमान \omterm{def:b3:topology:continuity}{संतत} है, और निष्कर्ष
        निकालिए: प्रत्येक $\varepsilon > 0$ के लिए ऐसा $n_0$ है कि $n \geq n_0$ और
        प्रत्येक अजालिका $t$ के लिए $\norm{\Delta_n(t)} \leq \varepsilon$। \emph{($(t_k,
        t_{k+1})$ पर
        $\varphi_n'(t) = f(t_k, \varphi_n(t_k))$, और $(t, \varphi_n(t))$ $(t_k, \varphi_n(t_k))$ से $(1 + M)\,T/n$ दूरी के भीतर है।)}
  \item समाकल रूप स्थापित कीजिए: सभी $t$ के लिए
        \[
        \varphi_n(t) = x_0 + \int_{t_0}^{t}
        f\bigl(s, \varphi_n(s)\bigr)\dd s +
        \int_{t_0}^t \Delta_n(s)\,\dd s,
        \]
        जिसमें बीच का समाकल्य खंडशः \omterm{def:b3:topology:continuity}{संतत} है।
  \item $(n_j)$ के अनुदिश सीमा लीजिए: दर्शाइए कि
        $f(s, \varphi_{n_j}(s)) \to f(s, \varphi(s))$
        \emph{एकसमान} रूप से (फिर $f$ का एकसमान \omterm{def:b3:topology:continuity}{सांतत्य}), और
        निष्कर्ष निकालिए
        \[
        \varphi(t) = x_0 + \int_{t_0}^t
        f\bigl(s, \varphi(s)\bigr)\dd s .
        \]
  \item निष्कर्ष निकालिए कि $\varphi$ $[t_0, t_0 +
        T]$ पर $\mathcal C^1$ है और $x' = f(t, x)$,
        $x(t_0) = x_0$ को हल कर देता है; रचना को $[t_0 - T, t_0]$ तक विस्तारित कीजिए
        (समय का उत्क्रमण)। \emph{यही पेआनो की प्रमेय है।}
\end{enumerate}

\medskip
\noindent\textbf{भाग III --- अद्वितीयता वास्तव में विफल हो जाती
है।} $\R$ पर $x' = 2\sqrt{\abs x}$, $x(0) = 0$ पर विचार कीजिए।
\begin{enumerate}[resume]
  \item जाँचिए कि $f(x) = 2\sqrt{\abs x}$ \omterm{def:b3:topology:continuity}{संतत} है पर $0$ के किसी भी \omterm{def:b3:topology:topology}{प्रतिवेश} पर
        लिप्शिट्ज़ नहीं।
  \item सत्यापित कीजिए कि $x \equiv 0$ और प्रत्येक $c \geq 0$ के लिए
        \[
        x_c(t) = \begin{cases} 0 & t \leq c,\\ (t - c)^2 & t
        \geq c,\end{cases}
        \]
        सभी $(0,0)$ से होकर जाने वाले हल हैं: भिन्न-भिन्न हलों का
        एक सातत्य।
  \item इस $f$ के लिए पिकार पुनरावृत्ति का तर्क (बानाख स्थिर
        बिंदु) कहाँ टूट जाता है?
\end{enumerate}

\medskip
\noindent\textbf{भाग IV --- विधि की सीमाएँ।}
\begin{enumerate}[resume]
  \item दर्शाइए कि पेआनो की प्रमेय अनंत विमा में विफल हो जाती
        है: हम शून्य अनुक्रमों की समष्टि $c_0$ में द्यूदोने का
        शास्त्रीय उदाहरण स्वीकार कर लेते हैं (आप उसे मान भी सकते
        हैं); इसके स्थान पर वह परिमित-विमीय अवयव सिद्ध कीजिए जो
        वहाँ विफल हो जाता है: $c_0$ (उच्चतम मानक) का संवृत इकाई
        गोला \omterm{def:b3:topology:compact}{संहत} नहीं है --- ऐसा परिबद्ध अनुक्रम प्रस्तुत कीजिए
        जिसका कोई अभिसारी उपानुक्रम न हो, और समझाइए कि भाग I का
        कौन-सा चरण टूट जाता है।
  \item सार लिखिए: कौन-सी परिकल्पनाएँ अस्तित्व देती हैं? और
        कौन-सी अस्तित्व तथा अद्वितीयता? दोनों प्रमेयों (पेआनो;
        कोशी--लिप्शिट्ज़) का ठीक-ठीक कथन साथ-साथ दीजिए।
\end{enumerate}

\medskip
\noindent\textbf{भाग V --- ऑस्गुड: लिप्शिट्ज़ से आगे
अद्वितीयता।} मान लीजिए $\omega\colon\intoo0\infty\to\intoo0\infty$ \omterm{def:b3:topology:continuity}{संतत}, अनवरोही है, और
\[
\int_0^1\frac{\dd r}{\omega(r)} = +\infty
\qquad(0 \text{ पर अपसरण}),
\]
तथा मान लीजिए $f$ $R$ पर $\norm{f(t, x) - f(t, y)} \leq
\omega\bigl(\norm{x - y}\bigr)$ को संतुष्ट करता है।
\begin{enumerate}[resume]
  \item जाँचिए कि $\omega(r) = Lr$ योग्य है (लिप्शिट्ज़), कि $\omega(r) = r\log\frac1r$
        (छोटे $r$ के लिए संततता से विस्तारित) योग्य है यद्यपि
        वह $O(r)$ \emph{नहीं} है, और यह कि $\omega(r) = 2\sqrt r$ योग्य नहीं है।
        प्रत्येक स्थिति में समाकल परिकलित कीजिए।
  \item मान लीजिए $x_1, x_2$ $x_1(t_0) = x_2(t_0)$ के साथ $[t_0, t_0 + T]$ पर समीकरण को हल करते
        हैं, और $\delta(t) =
        \norm{x_1(t) - x_2(t)}$। केवल समाकल रूपों से दर्शाइए कि $t_0 \leq s \leq t$ के
        लिए:
        \[
        \delta(t) \leq \delta(s) +
        \int_s^t\omega\bigl(\delta(v)\bigr)\dd v .
        \]
  \item (ऑस्गुड की प्रमेय) मान लीजिए किसी $t_1$ के लिए $\delta(t_1) > 0$, और
        मान लीजिए $\tau = \sup\{t \leq t_1 : \delta(t) =
        0\}$। $s \in \intoo\tau{t_1}$ के लिए $u(t) = \delta(s)
        + \int_s^t\omega(\delta(v))\,\dd v$ रखिए। दर्शाइए कि
        $\delta \leq
        u$, $u' = \omega(\delta) \leq \omega(u)$, और निष्कर्ष निकालिए
        \[
        \int_{\delta(s)}^{u(t_1)}\frac{\dd r}{\omega(r)}
        \leq t_1 - s \leq t_1 - \tau .
        \]
        अब $s \downarrow \tau$ लीजिए और समाकल के अपसरण से विरोधाभास निकालिए।
        निष्कर्ष निकालिए: \emph{किसी उभयनिष्ठ आरंभिक शर्त से
        होकर जाने वाले हल एक ही होते हैं} --- अर्थात् ऑस्गुड की
        शर्त के अंतर्गत अद्वितीयता।
  \item परिणाम निकालिए: कोशी--लिप्शिट्ज़ अद्वितीयता $\omega(r) = Lr$ की
        स्थिति है; समीकरण $x' =
        x\log\frac1{\abs x}$ ($0$ पर $0$ से विस्तारित) के
        हल अद्वितीय हैं, यद्यपि उसका दायाँ पक्ष $0$ पर
        लिप्शिट्ज़ नहीं है; और $x' = 2\sqrt{\abs x}$ के लिए $\int_0\frac{\dd r}{2\sqrt r}$ का अभिसरण ही
        वह बात है जो किसी हल को परिमित समय में $0$ छोड़ने देती
        है --- समाकल $\int_0^{h^2}\frac{\dd r}{2\sqrt r} = h$ के मान का $x_c$ के पलायन व्यवहार से
        मिलान कीजिए।
\end{enumerate}

\medskip
\noindent\textbf{भाग VI --- दरें, योजनाएँ, कीपें।}
\begin{enumerate}[resume]
  \item (समाकल ग्रोनवाल प्रमेयिका) मान लीजिए $e, \eta \geq 0$ $[t_0, t_0+T]$ पर \omterm{def:b3:topology:continuity}{संतत}
        है और सभी $t$ के लिए $e(t) \leq \int_{t_0}
        ^t\bigl(L\,e(s) + \eta(s)\bigr)\dd s$। दर्शाइए
        \[
        e(t) \leq \frac{\sup\eta}{L}\bigl(\eu^{L(t - t_0)} -
        1\bigr)
        \]
        \emph{($G(t) = \int_{t_0}^t(Le + \eta)$ रखिए, $G'
        \leq LG + \sup\eta$ पर ध्यान दीजिए, और $\eu^{-Lt}G(t)$ का
        अवकलन कीजिए।)}
  \item (ऑयलर एक दर के साथ अभिसरित होता है) अब मान लीजिए $f$
        $R$ पर $x$ में $L$-लिप्शिट्ज़ और $t$ में
        $L'$-लिप्शिट्ज़ है। प्रश्न 4 के दोष परिबंध को (मात्रात्मक
        बनाकर: $\norm{\Delta_n} \leq (L' + LM)T/n$) प्रश्न 17 के साथ मिलाकर सिद्ध कीजिए
        \[
        \sup_{[t_0, t_0+T]}\norm{\varphi_n - \varphi}
        \;\leq\; \frac{(L' + LM)\,T}{n}\cdot
        \frac{\eu^{LT} - 1}{L} = O\Bigl(\frac1n\Bigr),
        \]
        जहाँ $\varphi$ \emph{वही} हल है: लिप्शिट्ज़ आँकड़ों के साथ
        पूरा अनुक्रम अभिसरित होता है, और एक स्पष्ट दर के साथ ---
        किसी उपानुक्रम की आवश्यकता नहीं। इस परिकलन के बिना भी
        अद्वितीयता उपानुक्रमीय अभिसरण को \omterm{def:b3:complete:complete}{पूर्ण} अभिसरण तक क्यों
        उठा देती है?
  \item (योजना चुन लेती है) $x' = 2\sqrt{\abs x}$, $x(0) =
        0$ के लिए: दर्शाइए कि
        प्रत्येक ऑयलर बहुभुज सर्वथा शून्य है, अतः योजना हल $x \equiv 0$
        पर अभिसरित होती है; परंतु $x(0) = \varepsilon > 0$ से आरंभ करने पर वह (पहले
        $n \to \infty$, फिर $\varepsilon \to 0$) $t
        \mapsto t^2$ पर अभिसरित होती है, जो मूल बिंदु
        से होकर जाने वाला एक \emph{भिन्न} हल है। अनद्वितीयता
        सांख्यिक योजना की विक्षोभों के प्रति संवेदनशीलता के रूप
        में फिर उभर आती है।
  \item दर्शाइए कि $[t_0, t_0+T]$ पर $x' = f(t,x)$, $x(t_0) = x_0$ के \emph{समस्त} हलों का
        समुच्चय $\mathcal S$ (मान $\bar B(x_0, b)$ में) अरिक्त (भाग II), एकसमान
        रूप से $M$-लिप्शिट्ज़ और $\bigl(\mathcal C([t_0, t_0+T], \R^d),
        \norm\cdot_\infty\bigr)$ में संवृत है; आस्कोली से
        निष्कर्ष निकालिए कि $\mathcal S$ \emph{\omterm{def:b3:topology:compact}{संहत}} है। (क्नेज़र की
        प्रमेय यह भी जोड़ती है कि $\mathcal S$ \omterm{def:b3:topology:connected}{संबद्ध} है; हम उसे सिद्ध
        नहीं करेंगे।)
  \item उदाहरण पर क्नेज़र की परिघटना सत्यापित कीजिए: $[0, 1]$ पर
        $x' =
        2\sqrt{\abs x}$, $x(0) = 0$ के लिए दर्शाइए कि $x_\infty \equiv 0$ के साथ $\mathcal S = \{x_c : c \in \intcc0\infty\}$
        \emph{(किसी भी हल के लिए $c =
        \sup\{t : x(t) = 0\}$ लीजिए और $\{x > 0\}$ पर $(\sqrt x)' = 1$
        का समाकलन कीजिए)}, कि $c \mapsto x_c$ $\intcc0\infty$ (\omterm{def:b3:topology:locallycompact}{एक-बिंदु संहतीकरण},
        अर्थात् सीमा के रूप में जुड़ा $c = \infty$) से $\mathcal C([0,1])$ में \omterm{def:b3:topology:continuity}{संतत}
        है, और निष्कर्ष निकालिए कि $\mathcal S$ वस्तुतः \omterm{def:b3:topology:compact}{संहत} और \omterm{def:b3:topology:connected}{संबद्ध}
        है --- एक रेखाखंड के आकार की कीप।
  \item (पहुँच-योग्य समुच्चय) प्रश्न 20 से निष्कर्ष निकालिए कि
        प्रत्येक स्थिर $t$ के लिए \emph{पहुँच-योग्य समुच्चय}
        $\mathcal S(t) =
        \{x(t) : x \in \mathcal S\}$ \omterm{def:b3:topology:compact}{संहत} है; प्रश्न 21 के उदाहरण के लिए उसे परिकलित
        कीजिए और जाँचिए कि वह \omterm{def:b3:topology:connected}{संबद्ध} भी है: $\mathcal S(t) = \intcc0{t^2}$ --- अर्थात्
        प्रत्येक मध्यवर्ती अवस्था किसी न किसी हल से प्राप्त होती
        है।
\end{enumerate}

\medskip
\noindent\textbf{भाग VII --- पूरक: निर्भरता, इष्टतमता, और हाथ से
परिकलित एक योजना।}
\begin{enumerate}[resume]
  \item (\omterm{def:b3:topology:continuity}{संतत} निर्भरता) मान लीजिए $f$ $R$ पर $x$ में
        $L$-लिप्शिट्ज़ है, और मान लीजिए $x, y$ $t_0$ पर आरंभिक
        मानों $x_0, y_0$ वाले दो हल हैं। प्रश्न 17 की उपपत्ति को
        असमिका $e(t) \leq \norm{x_0
        - y_0} + \int_{t_0}^tL\,e(s)\dd s$ के अनुकूल बनाकर दर्शाइए
        \[
        \norm{x(t) - y(t)} \leq \norm{x_0 - y_0}\,
        \eu^{L(t - t_0)},
        \]
        और $x' = Lx$ पर जाँचिए कि यह परिबंध प्राप्त हो जाता है:
        ग्रोनवाल तीक्ष्ण है। इससे अद्वितीयता पुनः निकालिए
        ($x_0 =
        y_0$), और यह भी कि प्रवाह प्रतिचित्रण $x_0 \mapsto x(t; x_0)$ जहाँ भी
        परिभाषित है वहाँ नियतांक $\eu^{LT}$ के साथ लिप्शिट्ज़ है।
  \item (ऑस्गुड इष्टतम है) विलोमतः, मान लीजिए $\omega$ \omterm{def:b3:topology:continuity}{संतत},
        अनवरोही, $\intoo0\infty$ पर धनात्मक है, जिसके लिए $r \to 0^+$ पर $\omega(r) \to 0$,
        परंतु
        \[
        \int_0^1\frac{\dd r}{\omega(r)} < +\infty,
        \]
        और $f(x) = \omega(\abs x)$ को $f(0) = 0$ से विस्तारित कीजिए। दर्शाइए कि $\Omega(x) = \int_0^x\frac{\dd r}{\omega(r)}$
        $\intoc0{1}$ से $\intoc0{\Omega(1)}$ पर एक बढ़ता हुआ एकैकी आच्छादक प्रतिचित्रण
        है, कि उसका प्रतिलोम $g$ $g(0^+) = 0$ के साथ $g'
        = \omega(g)$ को हल करता
        है, और यह कि $g$, जिसे $t \leq 0$ के लिए $0$ से विस्तारित
        किया गया हो, $(0, 0)$ से होकर जाने वाला $x' = f(x)$ का ऐसा $\mathcal C^1$
        हल है जो $x \equiv
        0$ से भिन्न है \emph{($g'(0) = 0$ के लिए $\frac{g(t)}t$ को
        $\omega(g(t))$ से परिबद्ध कीजिए)}। निष्कर्ष निकालिए: प्रश्न 15 की
        अपसरण परिकल्पना कोई सुविधा नहीं, बल्कि अद्वितीयता की
        ठीक-ठीक सीमा है; $\omega(r) = 2\sqrt r$, $\Omega(x) = \sqrt x$ से भाग III पुनः प्राप्त
        कीजिए।
  \item (ऑयलर हाथ से परिकलित) $[0, 1]$ पर $x' = x$, $x(0) = 1$ के लिए:
        दर्शाइए कि ऑयलर बहुभुज $\varphi_n(1) = \bigl(1 + \frac1n\bigr)^n$ को संतुष्ट करता है। प्रसार
        \[
        \Bigl(1 + \frac1n\Bigr)^{n}
        = \eu\Bigl(1 - \frac1{2n} + O\Bigl(\frac1{n^2}\Bigr)
        \Bigr),
        \]
        सिद्ध कीजिए, अतः $t = 1$ पर त्रुटि $\frac{\eu}{2n} +
        O(n^{-2})$ है: यही प्रश्न 18
        की $O(\frac1n)$ दर है, ठीक नियतांक सहित। $n = 10$ के लिए संख्यात्मक
        जाँच कीजिए: $\eu \approx 2.71828$ के सामने $1.1^{10} = 2.59374$, अर्थात् त्रुटि $0.12454$,
        जिसकी तुलना $\frac{\eu}{20}
        \approx 0.13591$ से कीजिए।
\end{enumerate}
\end{problem}
